\section{Joint phase cancellation and primitive Dirichlet transversality}

The two marginal phase means used in NN-A036 can be replaced by a
proved mixed mean. This also answers the fixed primitive Dirichlet
analogue. Throughout this section a character and its conductor are
fixed as the height tends to infinity. No uniformity over conductors,
family average, full GUE conjecture, or simplicity assumption is used.
The character of conductor one denotes the Riemann zeta function.

\begin{theorem}[NN-A037: a joint arithmetic mode and definite normal motion]
Assume A. Let $L(s,\chi)$ be any fixed primitive Dirichlet $L$-function,
including $\zeta$. Write $a\in\{0,1\}$ for its parity and put
\[
 Q_\chi(s)=(q/\pi)^{(s+a)/2}\Gamma((s+a)/2),\qquad
 \Lambda_\chi(s)=Q_\chi(s)L(s,\chi),\qquad
 \Lambda_\chi(s)=\epsilon_\chi\overline{\Lambda_\chi(1-\bar s)}.
\]
Choose a square root of $\epsilon_\chi$ and the corresponding continuous
real phase $\theta_\chi$ for which
$Z_\chi(t)=e^{i\theta_\chi(t)}L(1/2+it,\chi)$ is real. Let
$N_\chi(T)$ count positive-ordinate nontrivial zeros with multiplicity,
and $N_{\chi,s}(T)$ count simple zeros. At a simple zero $\rho=1/2+i\gamma$
the specified norm shift has velocity
\[
 v_{\rho,\chi}=
 \frac{\rho L(\rho+1,\chi)}{L'(\rho,\chi)},\qquad
 r_{\rho,\chi}=\frac{|\Re v_{\rho,\chi}|}{|v_{\rho,\chi}|}.
\]
For every $\sigma>1$, let
$\phi_{\chi,\sigma}(t)=\arg L(\sigma+it,\chi)$ be the Euler-product
argument. Then
\begin{equation}\label{eq:A037-joint}
 \frac1{N_\chi(T)}\sum_{0<\gamma\le T}
 e^{2i\theta_\chi(\gamma)}
 \frac{L(\sigma+i\gamma,\chi)}
      {\overline{L(\sigma+i\gamma,\chi)}}\longrightarrow0.
\end{equation}
Every sum over zeros in this theorem retains multiplicity. In particular,
with $\delta_\gamma=\arctan(1/(2\gamma))$,
\begin{equation}\label{eq:A037-second}
 \frac1{N_\chi(T)}\sum_{0<\gamma\le T}
 \cos^2\bigl(\theta_\chi(\gamma)+
              \phi_{\chi,3/2}(\gamma)-\delta_\gamma\bigr)
 \longrightarrow\frac12.
\end{equation}
Consequently,
\[
 \liminf_{T\to\infty}\frac1{N_\chi(T)}
 \sum_{\substack{0<\gamma\le T\\\rho\ {
m simple}}}
 r_{\rho,\chi}^{\,2}\ge\frac16.
\]
At least a lower proportion $1/6$ of all zero occurrences are simple
and transverse; relative to the simple zeros the lower proportion is
$1/4$. More quantitatively, if $0<\eta<1/2$ and $G_{\chi,\eta}(T)$
counts simple zeros with $r_{\rho,\chi}\ge\eta$, then
\begin{equation}\label{eq:A037-counts}
 \liminf\frac{G_{\chi,\eta}(T)}{N_\chi(T)}
 \ge\frac{1-4\eta^2}{6(1-\eta^2)},\qquad
 \liminf\frac{G_{\chi,\eta}(T)}{N_{\chi,s}(T)}
 \ge\frac{1-4\eta^2}{4(1-\eta^2)}.
\end{equation}
For example, at $\eta=1/4$ these bounds are $2/15$ and $1/5$.
For each fixed $\varepsilon>0$, every zero counted at a fixed positive
threshold also satisfies
$|\Re v_{\rho,\chi}|\gg_{\chi,\eta,\varepsilon}
\gamma^{1-\varepsilon}$ as $\gamma\to\infty$.
No positive/negative sign proportion is asserted.
\end{theorem}

\begin{lemma}[Twisted first Hardy mode]\label{lem:A037-twist}
For every fixed $x>0$, under GRH for the specified primitive character,
\begin{align}\label{eq:A037-twist}
 \sum_{0<\gamma\le T}e^{2i\theta_\chi(\gamma)}x^{i\gamma}
 &=-\frac{i^a\epsilon_\chi^{-1}}{\sqrt{qx}}
    \sum_{n\le qxT/(2\pi)}\chi(n)\Lambda(n)e^{-2\pi i n/(qx)}
    +O_{\chi,x}(\sqrt T\log^2 T)\\
 &=O_{\chi,x}(T).\nonumber
\end{align}
Changing the convention at the upper endpoint changes the displayed
expression by an admissible error. For $q=1$, $x=1$, this is the
$\kappa=1$ formula of Arias de Reyna, Corollary~3.2, equation~(13).
The proof below gives the required twist and character extension,
including their constants; they are not inferred from that untwisted
statement.
\end{lemma}

\begin{proof}
Put
\[
 \Delta_\chi(s)=\epsilon_\chi Q_\chi(1-s)/Q_\chi(s).
\]
At a critical-line point,
$\Delta_\chi(1/2+it)^{-1}=e^{2i\theta_\chi(t)}$.
The gamma quotient is meromorphic with all of its zeros and poles real;
it is holomorphic and nonzero in the open upper half-plane. In every
fixed vertical strip, Stirling's formula gives, uniformly for $t\ge2\pi$,
\begin{equation}\label{eq:A037-stirling}
 \Delta_\chi(\sigma+it)^{-1}
 =\epsilon_\chi^{-1}(qt/(2\pi))^{\sigma-1/2}
 e^{i\{t\log(qt/(2\pi))-t+\pi(2a-1)/4\}}
 (1+O_\chi(t^{-1})).
\end{equation}
This follows directly by applying the usual Stirling expansion to
$\Gamma((s+a)/2)$ and $\Gamma((1-s+a)/2)$. The phase constant is
recorded because it supplies $i^a$ in \eqref{eq:A037-twist}.

First choose $T'$ in $(T,T+1)$ at distance $\gg_\chi1/\log T$ from
all zero ordinates. Such a choice is possible because the number of
zeros in this interval, or in a fixed enlargement of it, is
$O_\chi(\log T)$: see Iwaniec--Kowalski, Proposition~5.7(1),
equation~(5.27), pp.102--103. Their Theorem~5.6,
equations~(5.23)--(5.24), p.102, gives the Hadamard product and
partial-fraction formula. That formula gives
\[
 \frac{L'}L(\sigma+iT',\chi)=O_\chi(\log^2 T)
 \quad(-3/2\le\sigma\le 1+1/\log T).
\]
For clarity, this standard bound does not require simple zeros:
subtract the partial-fraction expansion at $2+iT'$; the poles at
distance at most one contribute at most $O(\log T)$ terms of size
$O(\log T)$, counting multiplicity. The remaining terms are bounded
by a constant times $\sum_{k\ge1}\log(T+k+2)/k^2=O(\log T)$.
The gamma and pole factors contribute $O(\log T)$. This proves the
claimed bound and also explains the admissible contour choice.

Integrate
\[
 \frac1{2\pi i}\frac{L'(s,\chi)}{L(s,\chi)}
          \Delta_\chi(s)^{-1}x^{s-1/2}
\]
around the rectangle with left side $\Re s=-3/2$, right side
$c=1+1/\log T'$, bottom at a fixed
$t_0\in(2\pi,2\pi+1)$ avoiding zero ordinates, and top at $T'$.
Any finite set of omitted lower zeros contributes $O_{\chi,x}(1)$.
There are no gamma poles or trivial zeros in this rectangle. The zeta
pole when $q=1$ is below it. The residue at a zero is its multiplicity
times $\Delta_\chi(\rho)^{-1}x^{\rho-1/2}$; GRH identifies this with
the summand in \eqref{eq:A037-twist}.

The bottom edge is $O_{\chi,x}(1)$. On the top edge the preceding
logarithmic-derivative bound and \eqref{eq:A037-stirling} give
$O_{\chi,x}(\sqrt T\log^2T)$ (in fact one logarithm can be saved).
On the left edge the functional equation and absolute convergence at
$5/2-it$ give $L'/L=O_\chi(\log(t+2))$, while
$|\Delta_\chi|^{-1}=O_\chi(t^{-2})$. Its integral is bounded.

Extending the right-edge integral down from $t_0$ to $2\pi$ costs
$O_{\chi,x}(1)$: on this fixed compact segment the integrand is
uniformly bounded as $c\downarrow1$. On the right edge use the absolutely convergent series
$L'/L=-\sum\chi(n)\Lambda(n)n^{-s}$ and Stirling. Its error is
\[
 O_{\chi,x}\left(\sum_{n\ge2}\frac{\Lambda(n)}{n^c}
                 \int_{2\pi}^{T'}t^{c-3/2}\,dt\right)
 =O_{\chi,x}(\sqrt T\log T).
\]
Here $T^{c-1}=O(1)$ and $\sum\Lambda(n)n^{-c}=O(\log T)$.
For the remaining integral set $u=t/(2\pi)$, $X=T'/(2\pi)$,
and $b=qx$. Its value is
\begin{equation}\label{eq:A037-stationary-integral}
 -\epsilon_\chi^{-1}e^{i\pi(2a-1)/4}
 b^{c-1/2}\sum_{n\ge2}\frac{\chi(n)\Lambda(n)}{n^c}
 \int_1^X u^{c-1/2}
 e^{2\pi i u\log(bu/(en))}\,du.
\end{equation}

Here are the full summation bounds for stationary phase. Set $r=n/b$.
Split $[1,X]$ into its intersection with $[r/2,2r]$ and the two
remaining intervals. On either remaining interval the derivative of
$u\log(u/(er))$ has absolute value at least $\log2$. The first
derivative test bounds each such integral by $O(X^{c-1/2})$.
Summing gives $O_{\chi,x}(\sqrt T\log T)$. The finitely many $r<1$
can be handled by the same test away from a bounded interval and add
at most this error.

On $[A,B]=[1,X]\cap[r/2,2r]$, the stationary-phase lemma of Arias
de Reyna, Lemma~2.5 (the underlying primary lemmas are Levinson and
Gonek), states that the integral is
\[
 {\bf1}_{A\le r\le B}\,r^c e^{i\pi/4}e^{-2\pi i r}
 +O\left(r^{c-1/2}
 +\frac{r^{c+1/2}}{|A-r|+\sqrt r}
 +\frac{r^{c+1/2}}{|B-r|+\sqrt r}\right).
\]
Its constants are uniform for $c-1/2$ in the compact interval in use.
Multiplication by the coefficient in
\eqref{eq:A037-stationary-integral} turns the main term into
\[
 -i^a\epsilon_\chi^{-1}b^{-1/2}
       \chi(n)\Lambda(n)e^{-2\pi i n/b}.
\]
The first remainder sums to $O_{\chi,x}(\sqrt T)$, since
$\sum_{n\le C T}\Lambda(n)n^{-1/2}\ll\sqrt T$.
Except for finitely many $n$, $A=r/2$, so the lower-endpoint
remainder has the same bound. If $r\le X/2$, then $B=2r$ and the
upper-endpoint remainder again has that bound. For $X/2<r\le2X$
it is at most
\[
 C_{\chi,x}\sqrt T\log(T+2)
 \sum_{bX/2<n\le2bX}
       \frac1{|n-bX|+\sqrt T}
 \ll_{\chi,x}\sqrt T\log^2T.
\]
The last bound follows by splitting the distances into unit intervals,
or comparing with $\int_0^{C T}(u+\sqrt T)^{-1}\,du$; it is valid
even when $bX$ is an integer. The constant number of small $n$ has
already been absorbed. This proves the displayed asymptotic at $T'$.
Chebyshev's bound $\sum_{n\le y}\Lambda(n)\ll y$
(Iwaniec--Kowalski, Section~2.2, equation~(2.12), p.33) gives its
$O_{\chi,x}(T)$ consequence without requiring cancellation in the
additive prime sum.

Finally, under GRH each summand on the left has modulus one, and
there are $O_\chi(\log T)$ zeros between $T$ and $T'$. The change
in the prime sum has $O_{q,x}(1)$ terms, each $O(\log T)$.
Thus the formula holds at every real cutoff $T$, with the stated
endpoint convention and error.
\end{proof}

\begin{lemma}[The precise fixed-character simplicity input]
Under GRH for a fixed primitive Dirichlet $L$-function,
\[
 \liminf_{T\to\infty}\frac{N_{\chi,s}(T)}{N_\chi(T)}\ge\frac23.
\]
This is a classical restricted-pair-correlation consequence, used here
as background rather than counted as a new simplicity theorem.
\end{lemma}
\begin{proof}
We specify the imported theorem and the multiplicity calculation to
avoid confusing a family average with a result for a fixed character.
Rudnick--Sarnak, \emph{Zeros of principal $L$-functions and random
matrix theory}, Theorem~1.2, specializes at degree $m=1$, correlation
order two, to every fixed primitive Dirichlet character. Its hypothesis
(1.7) is automatic here (and the theorem explicitly permits $m\le3$).
Their zeros are indexed with multiplicity; equation~(1.4) sums over
distinct indices, not necessarily distinct ordinates. The test functions
$f(x_1,x_2)=g(x_1-x_2)$ have Fourier support in
$(u,-u)$ with $|u|<1$, exactly the permitted
$|\xi_1|+|\xi_2|<2$. With $K(u)=\sin(\pi u)/(\pi u)$, this gives
\[
 \frac1N\sum_{j\ne k}g(\widetilde\gamma_j-
                         \widetilde\gamma_k)
 \longrightarrow\int_{\mathbb R}g(u)(1-K(u)^2)\,du.
\]
Here the indices run through the first $N$ positive zeros, and their
normalization is that of the cited theorem. Equal ordinates have
exactly equal normalized ordinates, which is all that is needed below.

Take nonnegative Schwartz tests $g$ with $g(0)=1$ and Fourier support
strictly inside $(-1,1)$. Every zero of multiplicity $m\ge2$ contributes
$m(m-1)\ge m$ to the sum on the left. Hence the fraction of occurrences
belonging to nonsimple zeros is bounded above by the displayed integral
in the limit. The same statement holds at cutoffs $T$ that include each
whole multiplicity block.

For completeness the endpoint optimization uses only admitted
Schwartz tests. For $0<\beta<1$, approximate
$\beta^{-1}{\bf1}_{[-\beta/2,\beta/2]}$ in $L^2$ and $L^1$ by
smooth real even functions $b_j$ supported inside that interval with
$\int b_j=1$. Then $g_j=|\widehat b_j|^2$ are nonnegative Schwartz,
$g_j(0)=1$, have Fourier support in $[-\beta,\beta]$, and converge in
$L^1$ to $g_\beta(u)=(\sin(\pi\beta u)/(\pi\beta u))^2$.
Since $0\le1-K^2\le1$, their integrals converge to
\[
 \int g_\beta(1-K^2)
 =\frac1\beta-1+\frac\beta3.
\]
Indeed $\widehat g_\beta(u)=\beta^{-1}(1-|u|/\beta)_+$ and
$\widehat {K^2}(u)=(1-|u|)_+$, so their product integral is
$1-\beta/3$. First take $N\to\infty$, then the smooth approximation,
then $\beta\uparrow1$. The nonsimple fraction is at most $1/3$.
Central zeros and finitely many endpoint choices do not affect this
positive-height limit.
\end{proof}

\begin{proof}[Proof of NN-A037]
For $\sigma>1$ the absolutely convergent expansion is
\[
 \frac{L(\sigma+it,\chi)}{\overline{L(\sigma+it,\chi)}}
 =\sum_{n,m\ge1}
 \frac{\chi(n)\mu(m)\overline{\chi(m)}}{(nm)^\sigma}
       e^{it\log(m/n)}.
\]
Its absolute coefficient sum is at most $\zeta(\sigma)^2$.
For every fixed pair $m,n$, Lemma~\ref{lem:A037-twist}, with
$x=m/n$, says that the normalized mixed sum tends to zero because
$N_\chi(T)\sim T\log T/(2\pi)$. For the positive-height counting
convention this is the calculation in the proof of
Iwaniec--Kowalski, Theorem~5.8, equations~(5.34)--(5.36), pp.104--105;
Rudnick--Sarnak equation~(2.11) records the corresponding two-sided
count. All such normalized sums have
absolute value at most one under GRH. Dominated convergence therefore
proves \eqref{eq:A037-joint}, without any uniform estimate as $m/n$
varies. This proves one mixed Fourier mode; it does not assert
independence or uniform distribution of the full joint phase.

Multiplication by $e^{-2i\delta_\gamma}$ changes the mean by $o(1)$.
To see this without a height-uniform error claim, fix $Y$, discard the
finitely many zeros below $Y$, and bound the remaining error by
$\sup_{\gamma>Y}2|\delta_\gamma|$; then take $T\to\infty$ and
$Y\to\infty$. Taking real parts and using
$\cos^2 u=(1+\cos2u)/2$ proves \eqref{eq:A037-second}.

At a simple zero,
$L'(\rho,\chi)=-i e^{-i\theta_\chi(\gamma)}Z_\chi'(\gamma)$,
where $Z_\chi'(\gamma)$ is real and nonzero. Since
$\arg(i\rho)=\pi-\delta_\gamma$, the exact velocity formula gives
\[
 r_{\rho,\chi}
 =\big|\cos(\theta_\chi(\gamma)+
              \phi_{\chi,3/2}(\gamma)-\delta_\gamma)\big|.
\]
The choice of square root of the root number changes the phase by
$\pi$ and therefore changes neither this expression nor the theorem.

Put $p=N_{\chi,s}(T)/N_\chi(T)$ and
$g=G_{\chi,\eta}(T)/N_\chi(T)$. The preceding simplicity lemma
gives $p\ge2/3-o(1)$. The contribution of all nonsimple occurrences
to \eqref{eq:A037-second} is at most $1-p$, so the simple contribution
has lower limit at least $1/6$. Also
\[
 \frac12+o(1)
 \le(1-p)+\eta^2(p-g)+g
 =1-(1-\eta^2)p+(1-\eta^2)g.
\]
It follows that
$g\ge p-[2(1-\eta^2)]^{-1}+o(1)$, which gives the first inequality
in \eqref{eq:A037-counts}. Dividing by $p$, bounded below by
$2/3-o(1)$, gives the second. Letting $\eta\downarrow0$ proves the
two asserted nonzero-velocity proportions. The numerical fractions at
$\eta=1/4$ are exact rational substitutions.

Finally, the Euler product gives
$|L(3/2+i\gamma,\chi)|\ge1/\zeta(3/2)$. The exact classical input
for the fixed-character Lindelof bound is Iwaniec--Kowalski,
\emph{Analytic Number Theory}, Section~5.7, Corollary~5.20,
equation~(5.60), printed p.116, following Theorem~5.19 on that page.
The hypotheses are GRH and the Ramanujan--Petersson coefficient
bound; the latter holds trivially for a Dirichlet character, since
$|\chi(n)|\le1$. The analytic conductor for fixed $\chi$ is
$\asymp_\chi 1+|t|$. Cauchy's derivative formula on a small
fixed circle, using the functional equation on its left half, gives
$|L'(1/2+i\gamma,\chi)|\ll_{\chi,\varepsilon}\gamma^\varepsilon$.
Explicitly, choose the circle radius smaller than $\varepsilon/4$;
Lindelof on the right and the gamma ratio bound on the left make the
supremum $O(\gamma^\varepsilon)$ after reducing the auxiliary exponent.
Thus $|v_{\rho,\chi}|\gg\gamma^{1-\varepsilon}$, and multiplication
by the fixed margin $\eta$ proves the final assertion.
\end{proof}

\paragraph{Exact scope and the Dedekind obstruction.}
This theorem is about every fixed primitive degree-one function. It
strengthens NN-A036 for $\zeta$ and proves its requested Dirichlet
analogue. A general Dedekind zeta function is not primitive. Repeated
Artin factors can produce nonsimple zeros systematically, and the
degree-one first-mode calculation cannot be imported for its product
without a new estimate. Nor do \eqref{eq:A037-joint} and
\eqref{eq:A037-second} distinguish the two signs of $\Re v$: the
sign of the real derivative $Z_\chi'(\gamma)$ is lost upon squaring.
The index-parity correlation and the magnitude-tail obstruction
recorded in Q-A036 therefore remain genuine sign-density gates.

\paragraph{Algebraic-geometric and representation translation.}
A primitive character gives its actual rank-one Artin representation
on the finite etale cyclotomic cover of
$\operatorname{Spec}\mathbb Z[1/q]$. Its root number supplies the
half-phase that makes the completed critical-line germ Real, including
for complex characters. At a reduced analytic zero the specified
deformation has differential
\[
 dL_t=L'(\rho,\chi)\,ds-\rho L(\rho+1,\chi)\,dt,
\]
so its tangent section has the stated normalized real normal
component; the ambient divisor is counted with its actual local
lengths. SGA~1, Expose~I, section~3, Proposition~1 is the local
differentials dictionary. On the compact dual of $\mathbb Q_{>0}$,
the mixed-mode theorem annihilates the first Hardy character times
each fixed prime-circle character, and absolute Euler convergence
extends that annihilation to the displayed exterior ratio. This is
a precise Fourier statement, not a full product-Haar or Frobenius
equidistribution theorem for analytic zeros. Rudnick--Sarnak's
restricted correlation theorem controls the nonreduced occurrence
mass sufficiently to obtain the explicit reduced-locus density.
No shifted arithmetic family, pure motive, or complexity separation
is supplied by this dictionary.

\paragraph{Primary locators used.}
Arias de Reyna, \emph{On the distribution (mod~1) of the normalized
zeros of the Riemann zeta-function}, arXiv:1403.7126v1,
Lemma~2.5, Corollary~3.2 equation~(13), and Section~4.4 give the
stationary-phase input and the untwisted comparison.
Rudnick--Sarnak, \emph{Zeros of principal $L$-functions and random
matrix theory}, author manuscript PDF pp.1--5, especially the
multiplicity convention on p.1, definitions (1.2)--(1.4) on p.3,
and Theorems~1.1--1.2 on pp.4--5, give precisely the fixed-character
restricted correlation input. Their distinction between primitive
and product functions is on p.6. The manuscript is available at
\url{https://empslocal.ex.ac.uk/people/staff/mrwatkin/zeta/rudnick-sarnak.pdf}.
The completion and root number use DLMF~25.15.5--25.15.6; the
Stirling expansion is DLMF~5.11.3 with its remainder formulation
in Section~5.11(ii). The absolute-speed clause uses
Iwaniec--Kowalski, \emph{Analytic Number Theory} (AMS Colloquium
Publications~53, 2004), Section~5.7, Corollary~5.20, equation~(5.60),
p.116, with Theorem~5.19's GRH and Ramanujan hypotheses checked
as above. The same book supplies the auxiliary Chebyshev bound at
Section~2.2, equation~(2.12), p.33; Hadamard and the logarithmic
derivative at Theorem~5.6, equations~(5.23)--(5.24), p.102; the
local zero count at Proposition~5.7(1), equation~(5.27), pp.102--103;
and the positive-height count in Theorem~5.8's proof,
equations~(5.34)--(5.36), pp.104--105.
The printed text was checked in an online OCR transcription
of the primary book, rather than a page image; the publisher record
is \url{https://doi.org/10.1090/coll/053}.
The new mixed-mode and transversality deductions
are proved in full above; no historical novelty claim is made.
